% Template for post or presentation abstracts submitted to ILAS 2022
% 1. Fill in the presenter's information (name, affiliation, email
% address) and talk information (title of presentation, abstract).
% 2. Compile to make sure there are no errors.
% 3. Save the .tex file for your abstract with a distinctive name,
% ideally "FamilynameGivenname.tex".
% 4. Upload your .tex file to https://clr.ie/129557
%
% * Please keep your LaTeX commands simple, and avoid the use of
%    user-defined macros.
% * Include details of co-authors and funding acknowledgements after the abstract.
% * Upload only the LaTeX file (with .tex extension).
% * Questions: Contact Niall Madden (Niall.Madden@NUIGalway.ie)

\documentclass[a4paper]{amsart}
\usepackage{amssymb}
\usepackage{hyperref}
\pagestyle{empty}
\date{}

%% IGNORE THE NEXT 4 LINES
\makeatletter %
\def\labelenumi{\theenumi.}
\def\theenumi{\@arabic\c@enumi}
\makeatother
%% STOP IGNORING NOW

\begin{document}

\title{Sharp bounds on the least eigenvalue of a graph determined from edge clique partitions}
\author{Domingos M. Cardoso} %Presenting author only
\address{Center for Research and Development in Mathematics and Applications - CIDMA, University of Aveiro} % Affiliation of presenting author
\email{dcardoso@ua.pt} % Email address of presenting author only

\maketitle

% The abstract  ***no more than one page***

Sharp bounds on the least eigenvalue of an arbitrary graph are presented. Necessary and sufficient (just sufficient) conditions for the lower (upper) bound to be attained are deduced using edge clique partitions. As an application, we prove that the least eigenvalue of the $n$-Queens' graph $\mathcal{Q}(n)$ is equal to $-4$ for every $n \ge 4$ and it is also proven that the multiplicity of this eigenvalue is $(n-3)^2$. Additionally, some results on the edge clique partition graph parameters are obtained.
%
%Please, follow the instructions at \url{http://ilas2022.ie/abstracts/}
%for submitting your abstract. Make sure you upload the \texttt{.tex}
%file, and not any others.
%
%Try to keep your file as simple as possible, so that it is easily
%rendered in various formats. Give the file a name that readily
%identifies it with you, for example \texttt{OlgaTaussky.tex}.
%
%If you wish to include citations, you can do so like this:
%\cite{my_key_for_OT49}. Use a bibitem key that is likely be be unique to
%you (e.g., don't use the default Math Reviews key).
%Any uncited bibliography entries will also
%appear. If you have none, then delete that section.
%
% Coauthor details here (optional - delete if not applicable)
% and funding acknowledgment (optional - delete if not applicable)

\bigskip

\emph{This is joint work with In\^es Ser\^odio Costa (University of Aveiro) and Rui Duarte (University of Aveiro).
Supported by the Center for Research and Development in Mathematics and Applications (CIDMA) which is financed by national funds through Funda\c{c}\~{a}o para a Ci\^{e}ncia e a Tecnologia (FCT), Grant UIDB/04106/2020.}
%
%\begin{thebibliography}{1}
%\bibitem{my_key_for_OT49}
%Olga Taussky.
%\newblock A recurring theorem on determinants.
%\newblock {\em  Amer. Math. Monthly}  56:672–676,  (1949).
%
%\bibitem{another_unique_key}
%Olga Taussky and John Todd.
%\newblock Cholesky, Toeplitz and the triangular factorization of
%symmetric matrices.
%\newblock {\em  Numer. Algorithms}, 41(2):197--202, 2006.
%\end{thebibliography}


\end{document}


% Please leave the next bit alone
%%% Local Variables:
%%% mode: latex
%%% TeX-master: "../ILAS-2022"
%%% End:
